%! Quadratic Functions Revisited — Unit 2, Lesson 2.1 — Homework
\documentclass[10pt]{article}
\usepackage{precalculus-article}
\usepackage{precalculus-boxes}
\newcommand{\TallMath}[1]{$\displaystyle #1\rule[-1.4em]{0pt}{3.2em}$}

\begin{document}

\pageheader{Unit 2, Lesson 2.1}{Homework}

\vspace{0.12in}

\begin{enumerate}[itemsep=10pt]

\item Read each feature straight off the form given.

\begin{enumerate}[label=(\alph*), itemsep=6pt]
  \item $g(x) = (x+1)^2 - 9$.

        Vertex: ( \blank{1.2cm} , \blank{1.2cm} ) \quad Axis of symmetry:
        $x = $ \blank{1.4cm} \quad Concave \blank{1.6cm} .

  \item $k(x) = -2(x-4)(x+2)$. \quad Zeros: \blank{1.2cm} and \blank{1.2cm}

        Since $a = -2$, the parabola opens \blank{1.8cm} .
\end{enumerate}

\item Let $f(x) = x^2 + 2x$. Find the average rate of change over each
      interval.

\begin{enumerate}[label=(\alph*), itemsep=6pt]
  \item Over $[0, 3]$:

        \begin{work}
          f(0) &= 0 \\
          f(3) &= (3)^2 + 2(3) \\
               &= 15 \\
          \text{AROC} &= \frac{15 - 0}{3 - 0} \\
                      &= 5
        \end{work}

        AROC $= $ \blank{1.8cm}

  \item Over $[3, 5]$:

        \begin{work}
          f(5) &= (5)^2 + 2(5) \\
               &= 35 \\
          \text{AROC} &= \frac{35 - 15}{5 - 3} \\
                      &= 10
        \end{work}

        AROC $= $ \blank{1.8cm}

  \item The AROC grew from (a) to (b), so the graph of $f$ is concave
        \blank{1.6cm} .
\end{enumerate}

\boxguard[18]
\item Let $q(x) = x^2 + x$. Complete the table and both difference rows.

\smallskip
\renewcommand{\arraystretch}{1.5}
\begin{tabular}{|c|c|c|c|c|c|}
\hline
$x$ & 0 & 1 & 2 & 3 & 4 \\
\hline
$q(x)$ & \blank{1cm} & \blank{1cm} & \blank{1cm} & \blank{1cm} & \blank{1cm} \\
\hline
\end{tabular}
\smallskip

First differences: \blank{1cm} , \blank{1cm} , \blank{1cm} , \blank{1cm}

Second differences: \blank{1cm} , \blank{1cm} , \blank{1cm}

The inputs step by $1$, so the second difference equals $2a$. That gives
$a = $ \blank{1.2cm} , which matches the formula.

\boxguard[18]
\item Two functions are recorded at inputs that step by $2$.

\smallskip
\renewcommand{\arraystretch}{1.4}
\begin{tabular}{|c|c|c|c|c|c|}
\hline
$x$ & 0 & 2 & 4 & 6 & 8 \\
\hline
$f(x)$ & 3 & 11 & 19 & 27 & 35 \\
\hline
$g(x)$ & 1 & 5 & 17 & 37 & 65 \\
\hline
\end{tabular}
\smallskip

\begin{enumerate}[label=(\alph*), itemsep=4pt]
  \item First differences of $f$: \blank{1cm} , \blank{1cm} , \blank{1cm} ,
        \blank{1cm}
  \item First differences of $g$: \blank{1cm} , \blank{1cm} , \blank{1cm} ,
        \blank{1cm}

        Second differences of $g$: \blank{1cm} , \blank{1cm} , \blank{1cm}
  \item One of these is linear and the other is quadratic. Say which is
        which, and justify each choice with the numbers you just computed.

        \vspace{0.05in}
        \writelines{3}
\end{enumerate}

\item A quadratic function is recorded at inputs that step by $1$, and its
      second differences are constantly $6$. What is its leading
      coefficient $a$?

\begin{enumerate}[label=\Alph*., itemsep=3pt]
  \item $3$
  \item $6$
  \item $1.5$
  \item $12$
\end{enumerate}

\end{enumerate}

\vspace{0.1in}

\boxguard
\begin{extensionbox}
Let $c(x) = x^3$, with outputs $0,\ 1,\ 8,\ 27,\ 64$ at $x = 0,1,2,3,4$. Its
first differences are $1, 7, 19, 37$ and its second differences are
$6, 12, 18$. Are the second differences constant --- and is $c$ quadratic?
Now take one more level of differences. What do you find?

\vspace{0.05in}
\writelines{2}
\end{extensionbox}

\vspace{0.1in}

\boxguard
\begin{notesbox}{Preview of Lesson 2.2}
Today the difference table stopped at level two, and level two was enough to
identify a quadratic. Lesson 2.2, \textit{Polynomial Functions and Rates of
Change}, keeps going: what a polynomial's rates of change say about where its
graph turns, and how far down the table you have to dig to name its degree.
\end{notesbox}

\end{document}
